\section{CUDA GPU 并行编程：仅使用全局内存的向量归约}



\subsection{课程引言与重要性}
本章是向量归约（Vector Reduction）系列的第一部分。向量归约不仅是独立的应用程序，更是理解 GPU 优化技术的基石。许多复杂应用（如矩阵乘法）内部都包含规约操作。掌握本章内容对于理解后续的优化策略至关重要。

\begin{tcolorbox}[colback=blue!5!white,colframe=blue!75!black,title=前置知识回顾]
在之前的 \textbf{向量加法} 中，将向量 A 与 B 对应元素相加存入 C，线程数等于元素数，每个线程独立处理一个索引。作为练习，可以尝试用一半的线程实现向量加法（即每个线程处理两个元素）。但本章的重点是 \textbf{向量归约}，其复杂度远高于向量加法。
\end{tcolorbox}

\subsection{向量归约的核心概念}

\subsection{什么是向量归约}
向量归约是指将单个输入向量（Input Vector）中的所有元素累加，并将最终总和存储在特定位置（通常是索引 0 处）的过程。

\begin{itemize}
    \item \textbf{输入}：一个包含 $N$ 个元素的向量，例如 $\{v_0, v_1, \dots, v_{15}\}$。
    \item \textbf{输出}：所有元素的总和，存储在 $input[0]$ 中。
    \item \textbf{难点}：存在数据依赖，不能像向量加法那样简单并行，需要多步骤迭代。
\end{itemize}

\subsection{基于树的归约方法（Tree-based Approach）}
归约过程呈现树状分支结构，每一步都将待处理元素数量减半：
\begin{enumerate}
    \item \textbf{第一阶段}：16 个元素 $\to$ 8 个有效结果。
    \item \textbf{第二阶段}：8 个有效结果 $\to$ 4 个有效结果。
    \item \textbf{第三阶段}：4 个有效结果 $\to$ 2 个有效结果。
    \item \textbf{最终阶段}：2 个有效结果 $\to$ 1 个最终总和。
\end{enumerate}
这种逐步减少活跃线程和数据量的过程被称为“规约”（Reduction）。

\subsection{单 Block 内的归约逻辑推导}

\subsection{第一步：相邻元素相加}
假设向量有 16 个元素，启动 16 个线程（Block Size = 16）。
\begin{itemize}
    \item \textbf{线程 0}：执行 $input[0] += input[1]$
    \item \textbf{线程 2}：执行 $input[2] += input[3]$
    \item \textbf{线程 4}：执行 $input[4] += input[5]$
    \item \dots 以此类推
\end{itemize}
\textbf{关键观察}：只有偶数 ID 的线程（0, 2, 4...）在工作，奇数 ID 线程（1, 3, 5...）处于空闲状态。虽然这看似浪费资源，但这是为了统一代码逻辑以支持后续步骤所必须的代价。

\subsection{为什么需要重复执行？}
第一步结束后，只得到了相邻两元素的和，并未得到全向量和。必须对新生成的值重复执行相同的加法逻辑，但每次参与的线程数和操作的元素间距都会变化。

\subsection{步幅（Stride）的概念}
步幅是指当前步骤中，参与运算的两个元素之间的索引距离。
\begin{itemize}
    \item \textbf{第 1 步}：线程0处理索引0和1，距离为 $1-0=1$，故 $stride=1$。
    \item \textbf{第 2 步}：线程 0 处理索引 0 和 2（因为索引 1 已在第 1 步被累加到 0 中），距离为 $2-0=2$，故 $stride=2$。
    \item \textbf{第 3 步}：线程 0 处理索引 0 和 4，距离为 $4-0=4$，故 $stride=4$。
    \item \textbf{规律}：$stride$ 从 1 开始，每步乘以 2（$1 \to 2 \to 4 \to 8 \dots$）。
\end{itemize}

\subsection{代码实现与过滤器设计}

\subsection{基础公式与问题}
最朴素的加法指令为：
\begin{lstlisting}
input[index] += input[index + stride];
\end{lstlisting}
直接执行此指令会导致两个严重问题：
\begin{enumerate}
    \item \textbf{无效计算}：奇数线程或不符合当前步幅要求的线程会执行多余操作。虽然不影响最终结果（因为后续步骤不会读取这些位置），但浪费算力。
    \item \textbf{内存越界}：当 $index + stride \ge N$ 时，访问非法内存地址，导致程序崩溃或数据损坏。
\end{enumerate}

\subsection{双重过滤器（Filters）}
为确保正确性，必须在加法指令前添加两个条件判断：

\subsection{过滤器1：边界检查}
必须确保读取的第二个元素在合法范围内：
\begin{lstlisting}
if (index + stride < n)
\end{lstlisting}
注意：这里检查的是 $index + stride$ 而非 $index$，因为加法操作依赖于偏移后的元素。

\subsection{过滤器2：步幅对齐检查}
只有满足当前步幅倍数的线程才能执行操作。通用公式为：
\begin{lstlisting}
if (tid % (2 * stride) == 0)
\end{lstlisting}
\textbf{推导}：
\begin{itemize}
    \item 当 $stride=1$ 时，$tid \% 2 == 0$，选中 0, 2, 4...（偶数线程）。
    \item 当 $stride=2$ 时，$tid \% 4 == 0$，选中 0, 4, 8...（4的倍数线程）。
    \item 当 $stride=4$ 时，$tid \% 8 == 0$，选中 0, 8...（8的倍数线程）。
\end{itemize}

\subsection{完整的循环结构}
结合上述分析，单 Block 内的归约核函数核心逻辑如下：
\begin{lstlisting}[caption={带过滤器的归约循环}]
// tid 为线程在Block内的局部IDfor (unsigned int stride = 1; stride < blockDim.x; stride *= 2) {
    // 双重过滤：步幅对齐 + 边界安全
    if ((tid % (2 * stride) == 0) && (tid + stride < blockDim.x)) {
        input[tid] += input[tid + stride];
    }
    // 【关键】块内同步屏障
    __syncthreads(); 
}
\end{lstlisting}

\begin{tcolorbox}[colback=red!5!white,colframe=red!75!black,title=同步的重要性]
\textbf{\texttt{\_\_syncthreads()}} 是必不可少的。它确保当前步幅下所有线程都完成了写
入操作后，才进入下一步幅的计算。如果没有同步，某些线程可能会在邻居线程尚未更新数据时就读取旧值，导致计算结果完全错误。
\end{tcolorbox}

\subsection{多 Block 归约与全局同步挑战}

\subsection{部分和（Partial Sums）}
当向量大小超过单个 Block 的处理能力时，需将向量切分为多个片段，每个 Block 处理一个片段。每个 Block 归约后得到的结果称为“部分和”（Partial Sum）。例如：
\begin{itemize}
    \item Block 0 处理前 8 个元素 $\to$ 部分和 $S_0$
    \item Block 1 处理后 8 个元素 $\to$ 部分和 $S_1$
\end{itemize}

\subsection{CUDA 的同步限制}
\textbf{CUDA 不支持跨 Block 的全局同步。} 无法在一个核函数内部让所有 Block 停下来等待彼此完成。这是因为 GPU 硬件调度器可能不会同时调度所有 Block，强制全局同步会导致死锁或巨大的性能开销。

\subsection{解决方案：多核函数策略}
利用核函数启动之间的隐式全局同步机制：
\begin{enumerate}
    \item \textbf{核函数 1}：启动多个 Block，每个 Block 独立计算一段数据的部分和。
    \item \textbf{Host 端同步}：CPU 等待核函数 1 完全执行完毕（此时所有部分和已安全写入全局内存）。
    \item \textbf{数据重排}：将分散在各 Block 起始位置的部分和收集到连续内存区域。
    \item \textbf{核函数 2}：启动新核函数（通常只需 1 个 Block），对部分和向量进行最终归约。
\end{enumerate}

\subsection{部分和的收集代码}
在核函数 1 的末尾，需要将本 Block 的部分和移动到输入向量的前端，以便核函数 2 能连续读取：
\begin{lstlisting}
// 仅由Block内的0号线程执行
if (tid == 0) {
    // 将本 Block 归约结果（位于 blockIdx.x * blockDim.x）
    // 移动到向量的第 blockIdx.x 个位置
    input[blockIdx.x] = input[blockIdx.x * blockDim.x];
}
\end{lstlisting}
\textbf{注意}：核函数 2 处理的元素数量等于核函数 1 启动的 Block 数量，而非原始向量大小。

\subsection{完整 GPU 核函数代码}

\begin{lstlisting}[caption={仅使用全局内存的向量归约完整核函数}]
__global__ void reduce_in_place(float* input, int n) {
    unsigned int tid = threadIdx.x;
    unsigned int index = blockIdx.x * blockDim.x + threadIdx.x;
    
    // 1. 块内归约：生成部分和
    for (unsigned int stride = 1; stride < blockDim.x; stride *= 2) {
        // 过滤器：步幅对齐 & 边界检查
        // 注意：此处边界检查应针对全局索引或块内偏移
        // 简化写法假设n是blockDim.x的整数倍且仅做块内归约
        if ((tid % (2 * stride) == 0) && (tid + stride < blockDim.x)) {
            input[index] += input[index + stride];
        }
        __syncthreads(); // 块内同步
    }
    
    // 2. 收集部分和到向量头部
    if (tid == 0) {
        input[blockIdx.x] = input[index]; // index此时为blockIdx.x * blockDim.x
    }
\end{lstlisting}

\subsection{总结与后续展望}
\begin{itemize}
    \item \textbf{本章重点}：理解了归约的树形结构、步幅变化规律、双重过滤器的必要性、块内同步的作用，以及多Block场景下的全局同步解决方案。
    \item \textbf{性能局限}：本版本仅使用全局内存，存在严重的非合并访问（Non-coalesced Access）和高延迟问题。每一步都在读写全局内存，效率极低。
    \item \textbf{未来优化}：下一章将引入 \textbf{共享内存（Shared Memory）} 来减少全局内存访问，随后还会讲解 \textbf{Warp Shuffle} 指令、原子操作及更高级的归约模式。
    \item \textbf{学习建议}：如果第一次未能完全理解，请重新阅读本章并手动推演 16 个元素的归约过程。理解索引映射是本节最大的难点，也是最重要的收获。
\end{itemize}

